\section{The information in a complete Bellman support record}
\label{sec:fullsupport125}

A paid Bellman support asserts one inequality on the entire feasible action
domain. Retaining only its value and a recovered derivative interval loses
part of that assertion. With an uncertain supporting slope, the support
assertion must still hold on the whole domain. We give an exact finite
convex representation of this restriction. This advances the joint interpolation
calculation in Section~\ref{sec:joint124}, whose information class explicitly
excluded the raw supports. It does not identify a capability exclusive to
neural representations. Smooth convex interpolation has an established
literature \citep{THG2017Interpolation122}; the result below identifies the
additional constraints licensed by the original Bellman payload.
Optimal lower bounds based on collected first-order oracle information
also predate this construction \citep{FloreaNesterov2024Oracle125}.
Our contribution is the completion of the declared NBO payload on its
finite constrained action domain, a constructive sufficiency proof, and
its use in original-policy certificates. A general new theory of optimal
oracle bounds is not claimed.

Fix one state, original risk law and optimal-future objective
$q=Q_x$ on $[0,A]$, where $A>0$. Assume $q$ is continuously
differentiable, $\mu$-strongly convex, and has $L$-Lipschitz derivative,
$0<\mu\le L$. Retain the paid value and derivative intervals, and
the proved derivative cap $|q'|\le M$. For each original support record
$j$, the information also asserts that there is a single number
$p_j\in[p_j^-,p_j^+]$ such that
\begin{equation}\label{eq:globalpayload125}
 q(x)\ge S_j(x;p_j):=
 W_j+p_j(x-a_j)+\frac\mu2(x-a_j)^2
 \qquad(0\le x\le A).
\end{equation}
Its anchor, lower level and slope interval must come from the same
original calculation. A fitted derivative is not a substitute. The
quantifier in \eqref{eq:globalpayload125} is ``there exists one slope
valid at every action,'' not a new slope for each action. Exact slopes
are included by setting $p_j^-=p_j^+$. The original finite-risk support
record has this quantifier; its outward slope interval encloses the
single slope of that supporting objective.
For convex $h$ below, enforcing the entire robust pointwise lower envelope
is equivalent to existence of such a shared slope. The shared variable is
an exact convex representation of that global requirement, not an assertion
of additional information beyond the full robust obstacle.

Let $\mathcal Q_{\rm full}$ be exactly the nonempty class satisfying
these properties and all retained value and derivative intervals.
There is at least one finite value interval. Additional formulas of the
investment model are not included unless explicitly encoded. Set
$K=L-\mu$ and $h(x)=q(x)-\mu x^2/2$. A support becomes the affine
function
\begin{equation}\label{eq:affinesupport125}
 \ell_j(x;p_j)=W_j-p_ja_j+\mu a_j^2/2+(p_j-\mu a_j)x.
\end{equation}
Thus $h'$ is nondecreasing and $K$-Lipschitz, and $h\ge\ell_j$.

For $\ell\le0\le r$ define the bounded-domain conjugate
\begin{equation}\label{eq:boundedconjugate125}
 \Phi_K(v;\ell,r)=\max_{\ell\le t\le r}
              \{vt-Kt^2/2\}.
\end{equation}
When $K>0$, put $t_*={\rm proj}_{[\ell,r]}(v/K)$; then
$\Phi_K=vt_*-Kt_*^2/2$ and $\partial_v\Phi_K=t_*$.
When $K=0$, $\Phi_0=\max\{v\ell,vr\}$.
These formulas are exact rational calculations for rational inputs.

\subsection{A finite convex program that retains every support}

For trial actions $a,b$, form the sorted distinct nodes
$0=x_0<\cdots<x_n=A$ from the endpoints, every paid value and
derivative location, every support anchor, and $a,b$.
Write $f_i=h(x_i)$ and $g_i=h'(x_i)$. Transform all paid
value and derivative bounds accordingly; impose the original derivative
cap $|g_i+\mu x_i|\le M$ and finite value boxes obtained by integrating
that cap from any paid value interval. On each adjacent interval put
\[
 d_i=x_{i+1}-x_i,\qquad s_i=g_{i+1}-g_i,\qquad
 w_i=f_{i+1}-f_i-g_id_i.
\]
For $K>0$ retain the joint interpolation inequalities
\begin{equation}\label{eq:adjacentfull125}
 0\le s_i\le Kd_i,\qquad
 s_i^2\le2Kw_i,\qquad s_i^2\le2K(s_id_i-w_i).
\end{equation}
For $K=0$ use $s_i=w_i=0$. Include one shared variable $p_j$ in
its supplied interval for each support, and add, for every node and support,
\begin{equation}\label{eq:obstaclefull125}
 f_i-\ell_j(x_i;p_j)\ \ge\
 \Phi_K(p_j-\mu a_j-g_i;-x_i,A-x_i).
\end{equation}
Denote the resulting finite feasible set by $\mathcal P_{\rm full}(a,b)$.
It has a linear number of adjacent interpolation constraints and
$(n+1)m$ support inequalities for $m$ retained supports. Duplicate
observations are constraints at one node. Both domain endpoints are
present even when no measurement was bought there.

\begin{theorem}[Exact interpolation above complete Bellman supports]
\label{thm:fullinterpolation125}
A list $(f_i,g_i,p_j)$ belongs to $\mathcal P_{\rm full}(a,b)$ if
and only if there is a function in $\mathcal Q_{\rm full}$ realizing
the node values and derivatives, with these same $p_j$ satisfying
\eqref{eq:globalpayload125}. One such interpolant has a continuous
piecewise-affine derivative. The feasible set is nonempty, compact and
convex. Consequently the exact worst compatible original contrast is
\begin{equation}\label{eq:fullcontrast125}
 U_{\rm full}(a,b)=
 \max_{\mathcal P_{\rm full}(a,b)}
 \left\{f_b-f_a+\frac\mu2(b^2-a^2)\right\}
 =\max_{q\in\mathcal Q_{\rm full}}\{q(b)-q(a)\}.
\end{equation}
The extrema are attained. Removing the support inequalities enlarges
the information class, so its exact upper contrast cannot be smaller
than \eqref{eq:fullcontrast125}. The reversed optimization gives the
exact lower contrast. These are information bounds under the stated
interface; they are not oracle-complexity bounds for unrestricted access
to the original economic formulas.
\end{theorem}

The proof supplies a constructive sufficiency argument, rather than
assuming that nodewise lower bounds imply a global support. On an adjacent
interval the construction uses derivative slopes $K,0,K$ for $h$,
with their lengths fixed by the endpoint jets and the required area.
The extra term in \eqref{eq:obstaclefull125} is essential. At an
interior node a positive supporting-slope discrepancy requires a
positive value gap even when $f_i\ge\ell_j(x_i;p_j)$ alone holds.

\subsection{Verified contrasts without an exact optimizer}

Let $z$ collect $f,g,p$ in a valid finite rational box, and write all
constraints as $c_j(z)\le0$. The adjacent constraints are convex
quadratics. Each support constraint is convex because it is a convex
function of an affine argument minus an affine function. For $K>0$
it is continuously differentiable, including at its clipping points.
At $K=0$ choose any valid subgradient of its maximum of two affine
functions. Equality constraints may be represented by two inequalities.

For the minimization objective
$H=f_a-f_b-\mu(b^2-a^2)/2$, choose any rational point $z_0$,
nonnegative rational multipliers $y_j$, and valid subgradients
$v_j\in\partial c_j(z_0)$. Put
\begin{align}
 \mathcal L_0&=H(z_0)+\sum_jy_jc_j(z_0),\nonumber\\
 R&=\nabla H+\sum_jy_jv_j,\nonumber\\
 B&=\mathcal L_0+
 \sum_k\min\{R_k(\underline z_k-z_{0k}),
                         R_k(\overline z_k-z_{0k})\}.
 \label{eq:fullresidual125}
\end{align}
Then $q(b)-q(a)\le-B$ for every compatible original objective.
The formula requires neither primal feasibility nor stationarity of
$z_0$, nor strong duality. All terms are independently checkable by
rational arithmetic. A solver status alone is not a certificate.
An implementation intersects this upper bound with all earlier valid
same-objective bounds, so an unhelpful numerical dual proposal cannot
weaken its returned certificate. The arithmetic, failed proposals and
all retained supports remain charged.

For the unchanged implemented future $J=q+\eta$, $0\le\eta\le\delta$,
an upper contrast $u$ at a distinct realized action licenses a policy
change only when $u+\delta<0$. Absolute bounds on the same $q(b)$,
and a lower bound on its original optimum, must separately establish
the declared regret and objective-width requirements. At identical
actions the actual contrast is zero. No regularity of $\eta$ is
assumed, and the optimization of $q$ is not relabeled as optimization
of an arbitrary discontinuous $q+\eta$.

\paragraph{An exact decision consequence and its comparator.}
The supplement gives a rational example in which omitting a paid support
leaves the exact upper contrast for the action change $1\to0$ at
$-1/4$, whereas retaining it gives $-5/8$. The additional guaranteed
original gain is $3/8$ under the same disclosed payload; arbitrary future
error adds $\delta$ to either contrast. A cheap direct evaluation of the
support at the incumbent attains the stronger bound in this example.
It therefore demonstrates the economic loss from omitting support
information, and does not establish a computational advantage over a
conventional method that retains that information. A separate exact
counterexample shows why checking a support only at interpolation nodes
does not ensure its validity between nodes. The new theorem supplies
the finite exact program for all coupled records, while the present
original-economy experiment executes the separately specified acquisition
service. An algebraic information result is not relabeled as an executed
conic economic comparison.

\section{Paid action resolution and policy certificates}
\label{sec:paidresolution125}

The preceding theorem describes what a complete paid record permits.
It does not promise that the record resolves the decision or that the
finite program is economical to solve. The original-economy study therefore
also tests a declared service which buys new original-risk supports when
its inherited qualified policy remains unresolved. It uses the retained
rounded-midpoint contraction of Proposition~\ref{prop:midpoint123}, the
same original Bellman oracle, and the existing full-record bounds for
each actual acceptance. Midpoint localization is not introduced as a new
root-finding principle. The additions here connect its required support
precision, a complete work bound, the actor's initial economic certificate,
and the separate precision of the returned original policy.
The allocation of inexact-oracle accuracy is an established numerical
optimization problem \citep{DGN2014Inexact122,VanDesselGlineur2024}.
The work statement below is a closed-form solution of the particular
guaranteed-width planning problem specified here; its original-risk and
complete-policy premises remain part of the result.

An actor with original regret certificate
$J(a_0)-\min q\le K_0$ supplies the initial optimizer interval
\begin{equation}\label{eq:warmbracket125}
 I_0=[a_0-\sqrt{2K_0/\mu},a_0+\sqrt{2K_0/\mu}]\cap[0,A].
\end{equation}
An outward radius and intersection with previously certified derivative
brackets preserve validity. Every conventional actor with the same
certificate receives the same reduction. A small training loss without
that original-regret certificate does not shrink the interval.
The executed warm service uses its retained derivative-envelope bracket;
\eqref{eq:warmbracket125} is an additional valid initialization rule,
not an unrecorded change to that frozen implementation.

\begin{theorem}[A complete precision and work account for paid localization]
\label{thm:paidwork125}
Fix a relative rounding allowance $0\le\xi<1/2$ and
$\gamma=1/2+\xi<1$. Let $w_0$ bound the width of $I_0$, and choose
planned widths $w_{k+1}\in[\gamma w_k,w_k]$. Query an actually
feasible rounded midpoint of the current certified interval with rounding
error at most $\xi$ times its actual width. Suppose an oracle can return
an original support with level gap $e=U-W$ and slope width $v=p^+-p^-$
satisfying, for the requested next width $t=w_{k+1}$,
\begin{equation}\label{eq:jointprecision125}
 h=\chi t,\qquad
 \chi=\begin{cases}\min\{1/2-\xi,\mu/(4K)\},&K>0,\\
                         1/2-\xi,&K=0,\end{cases}
 \qquad e\le\frac{\mu\chi t^2}{4},\qquad v\le\frac{\mu t}{4}.
\end{equation}
Use a query only while the actual interval width exceeds $t$;
otherwise its planned contraction is already achieved. Then both
displacements $\pm h$ are feasible from the midpoint, and the finite
support conversion gives a derivative interval of width at most $\mu t$.
The new optimizer interval has width at most $t$.

Assume, as a separately verified oracle premise, that all computations
needed for one such support, including its unsuccessful precision levels,
cost at most $C_0+C_1t^{-p}$, with $C_0,C_1\ge0$ and $p>0$.
For a target $0<w_*<w_0$, put
\[
 N=\left\lceil\frac{\log(w_0/w_*)}{\log(1/\gamma)}\right\rceil,
 \qquad w_j=w_*\gamma^{-(N-j)},\quad 1\le j\le N.
\]
After at most $N$ completed precision-qualified probes, the optimizer
interval has width at most $w_*$. Their total work is bounded by
\begin{equation}\label{eq:geometricwork125}
 NC_0+C_1w_*^{-p}\frac{1-\gamma^{pN}}{1-\gamma^p}.
\end{equation}
Among all nonincreasing planned-width schedules with guaranteed
contraction floor $\gamma$, final width at most $w_*$ and this same
declared nonincreasing per-probe cost bound, this schedule minimizes the
sum of those bounds. It is an optimum for the stated certificate-planning
problem, not for all possible original-model algorithms or actual
shared-host elapsed times. Querying every probe at the terminal target
instead gives the bound $NC_0+NC_1w_*^{-p}$.

To return a policy, let a paid original-value interval at a feasible
$b$ be $[W_b,U_b]$, of width $e_b$, and let
$r_b=\max\{|b-\underline a|,|b-\overline a|\}$ for the final
optimizer interval. Then
\begin{align}
 |b-a_q^*|&\le r_b,\nonumber\\
 W_b\le J(b)&\le U_b+\delta,\nonumber\\
 0\le J(b)-\min q&\le\delta+Mr_b.
 \label{eq:terminalpolicy125}
\end{align}
If the optimizer is certified to be stationary, $Mr_b$ can additionally
be replaced by $Lr_b^2/2$. This regret bound does not require the value
width $e_b$: a certified original optimizer interval already controls
the economic action error. In particular, an interval strictly inside
$(0,A)$ certifies stationarity. A returned action must separately satisfy
$e_b+\delta\le\omega$ and the declared regret tolerance. Any fresh
value calculation needed for this final gate is charged in addition to
\eqref{eq:geometricwork125}, as are actor acquisition, the inherited
service and the certificate algebra. The theorem does not turn a narrow
optimizer interval into an unpaid absolute value certificate.
\end{theorem}

The proof uses a finite feasible radius in
\eqref{eq:jointprecision125}; it does not assume that an implemented
finite-radius conversion attains a real-valued infimum. The oracle-work
premise includes the required slope enclosure and original continuation
accuracy. A terminal-solver floor, a maximum risk resolution, an action
lattice floor or a resource cap can prevent it from holding. The executed
study records such cases as unresolved and retains the last qualified
policy. It neither fits a favorable power law to the realized times nor
claims \eqref{eq:geometricwork125} as an empirical complexity theorem
without verifying that premise.
The experiment executes its frozen finite integration ladder and records
achieved contractions. It does not claim to execute the ideal
cost-minimizing planned-width schedule in the theorem.

Every accepted replacement still requires its own strictly negative
upper bound on $J(b)-J(a_0)$ under the same future controller, as well
as the original regret and width gates. The study targets a stricter
regret tolerance as an additional outcome; a failure to attain that
target does not revoke the already qualified incumbent. Failed and
rejected queries remain in its complete service bill. These distinctions
connect information acquisition to an economic action rather than to an
internal fitting error.

